\scanpage{103}
% archive foliation: 102
\uline{III}) Conclusion, \inserted{$M^{\alpha}_{\mathcal C_0}(\overline{\mathbb Q})\simeq$} $M^{\alpha}_{\mathcal C_0}(\mathbb C)$ [$=M^{\alpha}_{\mathcal C_0\mathbb Q}(\mathbb C)\simeq\operatorname{Hom}(R^{\alpha}_{\mathcal C_0},\mathbb C)$, où $R^{\alpha}_{\mathcal C_0}$ est l'algèbre finie étale sur $\mathbb Q$ qui représente $M^{\alpha}_{\mathcal C_0\mathbb Q}$].
\margin{$\overline{\mathbb Q}=$ clôture alg. de $\mathbb Q$ dans $\mathbb C$}

On sait que c'est aussi l'ens $C^{\alpha}_{\mathcal C_0}$ des \inserted{classes d'iso} cartes \inserted{comb.} orientées qui sont de type $(\mathcal C_0,\alpha)$, i.e. Donc $\Gamma=\operatorname{Aut}(\overline{\mathbb Q})$ opère dans $\unsure{C}^{\alpha}_{\mathcal C_0}$ --- en fait \unsure{la} \deleted{l'opération}\inserted{connaissance} du $\Gamma$-ens \deleted{\unsure{i.e. op.}} $C^{\alpha}_{\mathcal C_0}$ signifie \deleted{\unsure{(à}} la conn. (à iso can. près) du $\mathbb Q$-schéma $M^{\alpha}_{\mathcal C_0\mathbb Q}$ (ou du $\mathbb Z[\frac1{\delta!}]$-schéma $(M^{\alpha}_{\mathcal C_0})_{\mathbb Z[\frac1{\delta!}]}$ \ldots). Bien entendu, $\Gamma$ opère de façon compatible \inserted{(i.e. commutant)} avec les opérations de $\gamma$ sur $C^{\alpha}_{\mathcal C_0}$.

Rien qu'en regardant l'élém. $\tau$ de $\Gamma$, $\tau z=\bar z$, on voit que l'opération de $\Gamma$ sur $C^{\alpha}_{\mathcal C_0}/\gamma$ peut ne pas être triviale\ldots

On va essayer de déterminer l'opération de $\gamma\times\Gamma$ sur $C^{\alpha}_{\mathcal C_0}$ dans certains cas élémentaires (cartes sphériques\ldots)
